\documentclass[10pt,a4paper]{amsart}
\usepackage[margin=19mm]{geometry}
\usepackage{amsmath,amssymb,mathtools}
\usepackage[scr=rsfs,cal=euler]{mathalfa}
\usepackage{xfrac}
\usepackage[unicode,hidelinks,bookmarks=false]{hyperref}
\newcommand{\ie}{i.e.\ }
\newcommand{\cf}{cf.\ }
\newcommand{\resp}{respectively\ }
\newcommand{\Subsection}{\subsection}
\providecommand{\nobreakdash}{\nobreak\hbox{-}}
\setlength{\emergencystretch}{2em}
\multlinegap0pt
\usepackage{amssymb}
\usepackage{mathtools}
\usepackage[shortlabels]{enumitem}
\usepackage{tikz}
\usepackage{float}
\newtheorem{lemma}{Lemma}
\title{Minimal-Volume Equiangular Hyperbolic $4$-Polytopes}
\author{Andrey Egorov}
\date{Source: arXiv:2609.04258v1 (CC BY 4.0)}
\begin{document}
\maketitle

\noindent\emph{Source and licence.} Minimal-Volume Equiangular Hyperbolic $4$-Polytopes. Version arXiv:2609.04258v1 (CC BY 4.0), distributed by arXiv under the Creative Commons Attribution 4.0 International licence, http://creativecommons.org/licenses/by/4.0/, as recorded in the arXiv licence metadata for this exact version and shown on the abstract page https://arxiv.org/abs/2609.04258v1. Full licence notice: \texttt{licenses/arxiv-2609.04258-CC-BY-4.0.txt}.\par\smallskip

\noindent\textit{Editorial scope.} Counted unit: the article's lemma lem:coefficients (source0.tex lines 573--586). The article's complete original proof of it is delivered with it (source0.tex lines 588--622, one \texttt{proof} environment of 35 lines). No mathematics is written, repaired or re-derived: the statement and the proof are the article's own lines, in the article's own notation, and no retained line is substituted or re-flowed.  Retained uncounted as the source-local context the proof needs: lines 484--486; lines 494--496; lines 498--500; lines 548--552; lines 560--562.  Nothing was omitted from any retained span, so the package carries no omission record.  No retained line contains a citation, so the wrapper carries no bibliography and no bibliography entry.  No retained line contains a figure environment, an image-inclusion command or drawing code, so no image is delivered and the package declares no asset dependency.  Editorial formatting only, in addition: an AMS \texttt{amsart} preamble that restates the article's own declarations and macros used by the retained lines, namely the environments \texttt{lemma} on separate counters, together with the standard packages those lines need.  The retained environments print in this document as Lemma 1 in order of appearance, whereas the article prints the same items with the numbering of its own full document.  The complete native source of the article as distributed in the arXiv e-print for this exact version is bundled unchanged as \texttt{sources/209354/source0.tex}.
\begin{equation}\label{eq:beta-alpha1}
 \beta(\alpha_1)=\frac{\pi}{3}.
\end{equation}

\begin{equation}\label{eq:omega-derivative}
 \omega'(\alpha)=\frac{3\beta(\alpha)}{2\pi^2}.
\end{equation}

\begin{equation}\label{eq:omega-endpoint}
 \omega\left(\frac{\pi}{2}\right)=\frac1{16}.
\end{equation}

\begin{equation}\label{eq:a-b-def}
 a(\alpha)=\omega(\alpha)+\frac12-\frac{\alpha}{\pi},
 \qquad
 b(\alpha)=\frac12-\frac{\alpha}{2\pi}.
\end{equation}

\begin{equation}\label{eq:volume-simplex}
 v(\Delta^4_\alpha)=1+5a(\alpha)-5b(\alpha),
\end{equation}

\begin{lemma}\label{lem:coefficients}
For $\alpha_0\leq\alpha<\pi/2$, one has
\begin{equation}\label{eq:a-positive}
 a(\alpha)>0
\end{equation}
and
\begin{equation}\label{eq:4a-b-positive}
 4a(\alpha)-b(\alpha)>0.
\end{equation}
Moreover, for $\alpha_0\leq\alpha\leq\alpha_1$,
\begin{equation}\label{eq:11a-3b-positive}
 11a(\alpha)-3b(\alpha)>0.
\end{equation}
\end{lemma}

\begin{proof}
By \eqref{eq:omega-derivative},
\[
 a'(\alpha)=\frac{3\beta(\alpha)}{2\pi^2}-\frac1\pi<0
\]
for $\alpha<\pi/2$.  Equations \eqref{eq:omega-endpoint} and
\eqref{eq:a-b-def} give $a(\pi/2)=1/16$, proving
\eqref{eq:a-positive}.

Put $h=4a-b$.  Then
\[
 h'(\alpha)=\frac{6\beta(\alpha)}{\pi^2}-\frac{7}{2\pi}<0,
\]
because $\beta(\alpha)<\pi/2$ gives
$h'(\alpha)<-1/(2\pi)$.  On the other hand,
$h(\pi/2)=4/16-1/4=0$.  Hence $h(\alpha)>0$ for
$\alpha<\pi/2$.

Finally, put $d=11a-3b$.  In the simplex range,
\eqref{eq:beta-alpha1} gives $\beta(\alpha)\leq\pi/3$, and therefore
\[
 d'(\alpha)=\frac{33\beta(\alpha)}{2\pi^2}
 -\frac{19}{2\pi}<0.
\]
Indeed, $\beta(\alpha)\leq\pi/3$ gives
$d'(\alpha)\leq-4/\pi$.
At $\alpha=\alpha_1$, the simplex is Euclidean-degenerate, so the continuous
extension of \eqref{eq:volume-simplex} is zero.  Hence
\[
 d(\alpha_1)=v(C^4_{\alpha_1})-v(\Delta^4_{\alpha_1})
 =v(C^4_{\alpha_1})>0.
\]
Since $d$ is decreasing, this proves \eqref{eq:11a-3b-positive} on the whole
interval $[\alpha_0,\alpha_1]$.
\end{proof}

\end{document}
